\section{Scrum Architecture}
\label{sec:scrum}

Group-07 contributes the environment sector to the wider MODELBD Scrum of
Scrums. The architecture separates work carried out inside the group from
cross-group coordination and cohort-level review. It records only duties that
are supported by the project evidence and does not assign unverified titles.

\subsection{Scrum-of-Scrums Architecture}
\label{subsec:scrum-architecture}

Figure~\ref{fig:scrum-architecture} shows the three confirmed coordination
levels and the direction in which milestones, dependencies and review findings
move.

\begin{figure}[H]
\centering
\begin{tikzpicture}[
  level/.style={draw=cublue,rounded corners,fill=cugrey,text width=9.8cm,
    minimum height=1.1cm,align=center},
  arr/.style={<->,>=Stealth,thick,draw=cuteal},
  node distance=0.7cm
]
\node[level] (l3) {\textbf{Level 3: cohort oversight}\\Shared milestones,
review expectations and integration direction};
\node[level,below=of l3] (l2) {\textbf{Level 2: cross-group coordination}\\
Intended exchange of sector progress, dependencies and review feedback};
\node[level,below=of l2] (l1) {\textbf{Level 1: Group-07 environment team}\\
Source work, ERD, normalization, SQLite implementation and verification};
\draw[arr] (l3)--(l2);
\draw[arr] (l2)--(l1);
\end{tikzpicture}
\caption{Scrum-of-Scrums structure followed by Group-07.}
\label{fig:scrum-architecture}
\end{figure}

\subsection{Roles and Duties at Each Level}
\label{subsec:scrum-duties}

Table~\ref{tab:scrum-duties} separates confirmed participants from the duties
that can be supported at each level. Individual Group-07 responsibilities are
listed in Table~\ref{tab:roles}.

\begin{table}[H]
\centering
\caption{Confirmed duties and intended coordination at each Scrum level}
\label{tab:scrum-duties}
\begin{tabular}{C{0.12\textwidth}L{0.28\textwidth}L{0.49\textwidth}}
\toprule
Level & Confirmed participants & Duties supported by project evidence \\
\midrule
1 & Group-07 environment team & Collect and verify sources, design the ERD,
normalize records, implement SQLite, test the database and report the sector. \\
2 & Cross-group coordination & Intended coordination function: sector progress,
integration dependencies and review feedback are exchanged between groups in
the wider structure. \\
3 & Cohort oversight & Set common milestones and reporting expectations and
review whether sector deliverables can contribute to the larger project. \\
\bottomrule
\end{tabular}
\end{table}

\subsection{Intra-Level Coordination}
\label{subsec:scrum-intra}

At Level 1, source evidence, extraction rules, normalization decisions, the ERD
and the SQLite database are reviewed as connected work products. A change in
record grain can alter a key; that key can alter the load order, competency SQL
and report explanation. Trello exposes task state while Drive, Excalidraw and
Git retain the material reviewed at each handoff.

\subsection{Inter-Level Coordination}
\label{subsec:scrum-inter}

Level 2 represents the intended cross-group coordination layer within the wider
Scrum-of-Scrums structure. For Group-07, the available project files mainly
show common milestones and review expectations rather than a direct exchange
of datasets, ERD components or implementation work with another group. Level 3
supplies the common academic schedule and review expectations. The report does
not attach names to these levels because the available evidence confirms the
coordination structure but not individual appointments.

\subsection{Data Collection and Processing}
\label{subsec:data-collection-processing}

Collection enters the database workflow only after an authority, scope and
record-grain review. Spreadsheets and CSV files can supply structured records;
PDF publications provide interpretation or manually checked values but are not
counted as database rows merely because they were collected. The review step
in Figure~\ref{fig:collection-processing} prevents an unverified source
from moving directly into normalization.

\begin{figure}[H]
\centering
\begin{tikzpicture}[
  procstep/.style={draw=cublue,rounded corners,fill=cugrey,text width=2.55cm,
    minimum height=1.1cm,align=center,font=\small},
  decision/.style={diamond,draw=cuteal,fill=white,aspect=2.0,align=center,
    font=\small},
  arr/.style={-{Stealth[length=2mm]},draw=cuteal,thick},
  node distance=0.55cm
]
\node[procstep] (discover) {Discover\\candidate source};
\node[procstep,right=of discover] (verify) {Verify\\authority and scope};
\node[decision,right=0.75cm of verify] (pass) {Suitable?};
\node[procstep,right=0.75cm of pass] (extract) {Extract\\named blocks};
\node[procstep,below=1.05cm of extract] (trace) {Attach\\source trace};
\node[procstep,left=0.8cm of trace] (approve) {Approve for\\normalization};
\node[procstep,below=1.05cm of verify] (record) {Record exclusion\\reason};
\node[procstep,left=0.8cm of record] (stop) {Retain in the\\source register};
\draw[arr] (discover)--(verify);
\draw[arr] (verify)--(pass);
\draw[arr] (pass)--node[above,font=\scriptsize]{yes}(extract);
\draw[arr] (extract)--(trace);
\draw[arr] (trace)--(approve);
\draw[arr] (pass.south) -- node[right,font=\scriptsize]{no} ++(0,-0.45) -|
  (record.north);
\draw[arr] (record)--(stop);
\end{tikzpicture}
\caption{Source review workflow before normalization.}
\label{fig:collection-processing}
\end{figure}

Of \mCellsRead{} inspected cells, \mCellsMissing{} are missing or unusable
(\mCellsMissingPct\%). This percentage describes source-cell quality, not rows
deleted from the final database. Sections~\ref{sec:data} and \ref{sec:etl}
explain the operational processing and its row accounting.
Table~\ref{tab:source-review-checks} states the questions used at this decision
point and the outcome recorded for each answer.

\begin{table}[H]
\centering
\caption{Checks applied at the source-review decision point}
\label{tab:source-review-checks}
\begin{tabular}{L{0.24\textwidth}L{0.31\textwidth}L{0.35\textwidth}}
\toprule
Check & Question & Recorded outcome \\
\midrule
Authority and scope & Is the publisher traceable and the subject relevant to
Bangladesh? & Continue or retain the exclusion reason. \\
Identifier and grain & Can records be identified at a supported station,
district, river or time grain? & Define the block grain or exclude it. \\
Structure & Can repeated observations be extracted without guessing headings
or values? & Name the block and preserve its source location. \\
Model fit & Does the material connect to the approved ERD? & Approve it for
normalization or keep it only as supporting material. \\
\bottomrule
\end{tabular}
\end{table}
